\section{从 CNN 到 ViT：如何把图片变成一串 Token？}
\label{intro:19b-vit}
\phantomsection\label{concept:vit}

上一章的卷积核可以在图片上\textbf{重复扫描局部窗口}。
但是文字 Transformer 不用卷积扫描每一个词，
也能通过 Attention 让不同位置交换信息。
一个自然的问题是：\textbf{能不能把图片也组织成一串向量}，
再用已经学过的 Attention 来处理？

\subsection{第一步：把一张图片切成 Patch}

先用一张只有一个通道的 $4\times4$ 小图，
不涉及颜色，也不用真实照片：
\[
I=\begin{bmatrix}
1&1&0&0\\
1&0&0&1\\
0&1&1&1\\
0&0&0&1
\end{bmatrix}.
\]
每个 Patch 是一个 $2\times2$ 方块，不重叠，
所以从左上到右下得到\textbf{四个 Patch}。
将每个 Patch 按行展开为长度 $4$ 的\textbf{行向量}，
排成一个矩阵：
\begin{equation}
X_{\rm patch}=
\begin{bmatrix}
1&1&1&0\\
0&0&0&1\\
0&1&0&0\\
1&1&0&1
\end{bmatrix}\in\R^{4\times4}.
\label{eq:vit-four-patches}
\end{equation}
第一条是原图左上角 $(1,1,1,0)$；
第四条来自右下角 $(1,1,0,1)$。
这里的四行指\textbf{四个位置}，
每行四个数字指\textbf{Patch 中的像素特征}，不要把两条轴混淆。

\subsection{第二步：给每个 Patch 学一条短向量}

文字的 Embedding 使用 Token ID 查表，
图像的 Patch 通常使用一套\textbf{共享的线性投影}：
\[
W_{\rm patch}=
\begin{bmatrix}
1&0&0\\
0&1&0\\
0&0&1\\
0.5&0.5&0
\end{bmatrix}\in\R^{4\times3}.
\]
按主线的行向量约定，右乘得到
\begin{equation}
\boxed{
E_{\rm patch}=X_{\rm patch}W_{\rm patch}
=\begin{bmatrix}
1&1&1\\
0.5&0.5&0\\
0&1&0\\
1.5&1.5&0
\end{bmatrix}\in\R^{4\times3}.
}
\label{eq:vit-patch-embedding}
\end{equation}
这是\textbf{人为设置的教学参数}，不代表预训练视觉模型的权重。
训练后，$W_{\rm patch}$ 可以学会将像素组合为有用的特征。
输入位置数由 Patch 数决定，隐藏维度由投影的列数决定。

\subsection{第三步：如果只给四行向量，怎么分清图像位置？}

回想本书第四部分的换位实验：
没有位置信息的 Self-Attention 对输入行置换是等变的。
因此 ViT 也需要把\textbf{空间位置}告诉模型。
例如把四行位置向量 $P_{\rm patch}\in\R^{4\times3}$
加到投影后：
\[
Z_0=E_{\rm patch}+P_{\rm patch}\in\R^{4\times3}.
\]
这里 $P_{\rm patch}$ 既可以来自可训练位置参数，
也可以使用其他位置机制，不强制等同于语言模型的逐词位置编码。

\subsection{第四步：Attention 如何读图像位置？}

对四个 Patch 的行向量进行 Q/K/V 投影，
普通\textbf{图像编码器}的双向 Self-Attention
可以让一块读取同一图片其他所有块：
\[
Q=Z_0W_Q,\qquad K=Z_0W_K,\qquad V=Z_0W_V.
\]
若每头维度是 $3$，
则 $QK^\top/\sqrt3$ 的形状为 $4\times4$，
Softmax 后仍是 $4\times4$，
输出 $AV$ 为 $4\times3$。
这里通常\textbf{不使用语言生成时的因果掩码}，
因为整张图已经可用；如果是图像生成或其他特殊任务，
则可能采用不同的注意力可见性规则。

有些图像分类 ViT 在 Patch 序列前加一个可训练的
\textbf{分类 Token}（常记作 \texttt{[CLS]}）。
本例原本是 $4\times3$，
加上这一个特殊位置就成为 $5\times3$，
注意力分数也从 $4\times4$ 变成 $5\times5$。
分类头可读取 CLS 输出，也有使用池化等其他方案。
\textbf{CLS 是一种设计选择，而不是每个视觉 Transformer 都必须采用。}

\subsection{与 CNN 的区别，以及现实尺寸}

CNN 的卷积核倾向于直接引入\textbf{局部性与权重共享}；
ViT 则把图片切成 Patch，再学习如何跨 Patch 组合表示。
并不能由“Attention 看得更远”直接推出 ViT
在任何图像任务、数据规模和硬件条件下都优于 CNN。

例如 $224\times224$ 彩色图片切成 $16\times16$ 的 Patch，
如果宽高恰好可整除，就有
$(224/16)^2=196$ 个 Patch；
每块展开后有 $16\times16\times3=768$ 个像素数。
若投影宽度也取 $768$，则 Patch 矩阵形状是
$196\times768$，投影矩阵形状是 $768\times768$。
这些是\textbf{尺寸说明}而非对某个 checkpoint 的性能承诺。
更多模型设计可参考 Dosovitskiy 等人的
\href{https://arxiv.org/abs/2010.11929}{An Image is Worth 16x16 Words}。

\begin{warningbox}
Patch 展平与词表 ID 查 Embedding 是不同的输入接口：
前者从像素块经线性投影产生连续特征，
后者从离散整数查可训练表。
它们进入 Transformer 后都可以按位置排成向量序列，
却不能据此断言所有视觉或语言系统有完全相同的训练目标。
\end{warningbox}

\par\medskip\noindent{\sffamily\footnotesize\color{blue}\textbf{读完自测}}\quad
\phantomsection\label{check:19b-vit}
本章的四个 Patch 投影后是什么形状？
如果添加一个 CLS Token，Attention 分数矩阵是什么形状？
\quad{\footnotesize\hyperref[answer:19b-vit]{核对解答。}}\par\medskip

\begin{conclusion}
从局部卷积到视觉 Transformer，改变的是如何构造和交换图像特征：
Patch 投影提供序列，位置表示帮助保留空间身份，
Attention 让不同 Patch 之间的信息能够交互。
\end{conclusion}
